\section{Appendix B: Basics of matrix analysis} \label{subsec:Matrix-analysis}

In this section, we review some basic facts in matrix analysis that are useful for the development in Chapter~\ref{{cha:matrix-perturbation}}. 
 

\subsection{Eigenvalues and eigenvectors}

We start with the definition of eigenvalues and eigenvectors of a real
square matrix.

\begin{definition}[Eigenpair]\label{def:eigenpair} Fix a square
matrix $\bm{A}\in\mathbb{R}^{n\times n}$. We say $(\lambda,\bm{v})\in\mathbb{C}\times\mathbb{C}^{n}$
is an \emph{eigenpair} of $\bm{A}$ if $\bm{v}$ is nonzero and 
\begin{equation}
\bm{A}\bm{v}=\lambda\bm{v}.\label{eq:def-eigenpair}
\end{equation}
In particular, the scalar $\lambda$ is an eigenvalue of $\bm{A}$
and $\bm{v}$ is the associated eigenvector. \end{definition}

In words, an eigenvector $\bm{v}$ of $\bm{A}$ does not change its
direction when multiplied by $\bm{A}$ and the eigenvalue $\lambda$
measures the amount of expansion or shrinkage in this linear transformation.
Several remarks regarding Definition~\ref{def:eigenpair} are in
order.
\begin{itemize}
\item The eigenpair (i.e.~the eigenvalue and the associated eigenvector) is only defined
for square matrices. This is implicitly assumed in the definition
of the eigenpair (see~(\ref{def:eigenpair})).
\item Every $n\times n$ square matrix has exactly $n$ (complex) eigenvalues,
possibly with repetitions. To see this, one can transform (\ref{def:eigenpair})
into the identity
\[
\left(\lambda\bm{I}_{n}-\bm{A}\right)\bm{v}=\bm{0}.
\]
Consequently, the scalar $\lambda$ is an eigenvalue of $\bm{A}$
if and only if 
\[
\mathsf{det}\left(\lambda\bm{I}_{n}-\bm{A}\right)=0,
\]
where $\mathsf{det}(\cdot)$ is the determinant of a matrix. Notice that the above relation specifies a polynomial equation of
degree $n$, which by the fundamental theorem of algebra, has exactly
$n$ roots in the complex field (and hence $n$ complex eigenvalues
of~$\bm{A}$).
\item A real matrix can have complex eigenvalues and
complex eigenvectors. Take for instance the matrix
\[
\bm{A}=\left[\begin{array}{cc}
3 & -2\\
4 & -1
\end{array}\right],
\]
where $1+2i$ and $1-2i$ are its two complex eigenvalues. Here $i$
denotes the imaginary unit.
\end{itemize}
With the last remark being said, there exists one important set of
real matrices --- \emph{real symmetric matrices} --- of which the eigenvalues
must be real. This claim is formalized in the following theorem.  

\begin{theorem}[Spectral theorem]\label{thm:spectral} Let $\bm{A}\in\mathbb{R}^{n\times n}$
be a real symmetric matrix. Then there exists $n$ real eigenvalues
$\lambda_{1},\cdots\lambda_{n}\in\mathbb{R}$ and an orthonormal basis
$\bm{u}_{1},\cdots\bm{u}_{n}\in\mathbb{R}^{n}$ such that $\bm{A}\bm{u}_{i}=\lambda_{i}\bm{u}_{i}$.
Equivalently, a real symmetric matrix $\bm{A}$ enjoys the following
spectral decomposition 
\begin{equation}
\bm{A}=\sum_{i=1}^{n}\lambda_{i}\bm{u}_{i}\bm{u}_{i}^{\top}.\label{eq:spectral-decomposition-summation-form}
\end{equation}
%where $\{\lambda_{i}\}_{1\leq i\leq n}$ are the real eigenvalues of $\bm{A}$ and $\{\bm{u}_{i}\}_{1\leq i\leq n}$ is an orthonormal basis of~$\mathbb{R}^{n}$.

\end{theorem}

Since the eigenvalues of real symmetric matrices are guaranteed to
be real numbers, we can therefore order them. From now on, denote
by $\lambda_{i}(\bm{A})$ the $i$th largest eigenvalue of a real
symmetric matrix $\bm{A}\in\mathbb{R}^{n\times n}$. In addition,
we define 
$$\lambda_{\max}(\bm{A})\coloneqq\lambda_{1}(\bm{A}) \quad \mbox{and}
\quad \lambda_{\min}(\bm{A})\coloneqq\lambda_{n}(\bm{A}). $$

To make the formula (\ref{eq:spectral-decomposition-summation-form})
more compact, we denote 
\[
\bm{U}=\left[\begin{array}{ccc}
\rvert &  & \rvert\\
\bm{u}_{1} & \cdots & \bm{u}_{n}\\
\rvert &  & \rvert
\end{array}\right]\in\mathbb{R}^{n\times n},\;\text{and}\;\bm{\Lambda}=\left[\begin{array}{cccc}
\lambda_{1} & 0 & \cdots & 0\\
0 & \lambda_{2} & \cdots & 0\\
\vdots & \vdots & \ddots & \vdots\\
0 & 0 & \cdots & \lambda_{n}
\end{array}\right]\in\mathbb{R}^{n\times n}.
\]
As a result, the decomposition (\ref{eq:spectral-decomposition-summation-form})
can be rewritten as 
\begin{equation}
\bm{A}=\bm{U}\bm{\Lambda}\bm{U}^{\top},\label{eq:spectral-decomposition-multiplication-form}
\end{equation}
in which a real symmetric matrix is decomposed into the multiplications
of orthonormal matrices (i.e.~$\bm{U}$) and a real diagonal matrix
(i.e.~$\bm{\Lambda}$). Another useful way to interpret spectral
decomposition is through the lens of diagonalization. Simple algebraic
manipulation reveals an equivalent form of (\ref{eq:spectral-decomposition-multiplication-form})
\begin{equation}
\bm{U}^{\top}\bm{A}\bm{U}=\bm{\Lambda}.\label{eq:spectral-decomposition-diagonalization-form}
\end{equation}
This essentially tells us that any real symmetric matrix $\bm{A}$
can be \emph{diagonalized }by an orthonormal matrix~(i.e.~$\bm{U}$).

\subsection{Singular values and singular vectors}

As we have pointed out in the remarks following Definition~\ref{def:eigenpair},
a rectangular matrix does \emph{not} have well-defined eigenvalues
and eigenvectors. It is then natural to ask whether there is a similar
decomposition of a general matrix to that of a real symmetric matrix
(cf.~(\ref{eq:spectral-decomposition-multiplication-form})). Fortunately,
a general matrix admits another useful representation: singular value
decomposition (SVD).

\begin{definition}[Singular value decomposition]\label{def:SVD}Let
$\bm{A}\in\mathbb{R}^{m\times n}$ be a general matrix. Without loss
of generality, assume that $m\leq n$. There exists a pair of orthonormal
matrices $\bm{U}\in\mathbb{R}^{m\times m}$ and $\bm{V}\in\mathbb{R}^{n\times n}$
such that 
\begin{equation}\label{eq:SVD}
\bm{A}=\bm{U}\bm{\Sigma}\bm{V}^{\top},
\end{equation}
where $\bm{\Sigma}$ is a rectangular diagonal matrix with 
\[
\bm{\Sigma}=\left[\begin{array}{cccccc}
\sigma_{1} & 0 & \cdots & \cdots & 0 & 0\\
0 & \sigma_{2} & \cdots & \cdots & 0 & 0\\
\vdots & \vdots & \ddots & \vdots & \vdots & \vdots\\
0 & 0 & \cdots & \sigma_{m} & 0 & 0
\end{array}\right]\in\mathbb{R}^{m\times n}
\]
and $\sigma_{1}\geq\sigma_{2}\geq\cdots\geq\sigma_{m}\geq0$. The
decomposition (\ref{eq:SVD}) is called singular value decomposition,
where the nonnegative real numbers $\{\sigma_{i}\}_{1\leq i\leq n}$
are the singular values of $\bm{A}$ and the columns in $\bm{U}$
(resp.~$\bm{V}$) are the left (resp.~right) singular vectors of
$\bm{A}$. \end{definition}

Similar to the eigenvalues of real symmetric matrices, the singular
values of a general matrix are guaranteed to be real. \yc{a bit redundant since the above already said all singular values are above zero (hence must be real).}As a result,
for any $1\leq i\leq\min\{m,n\}$, we define $\sigma_{i}(\bm{A})$
to be $i$th largest singular value of $\bm{A}$, define $\sigma_{\max}(\bm{A})\coloneqq\sigma_{1}(\bm{A})$
(resp.~$\sigma_{\min}(\bm{A})\coloneqq\sigma_{\min\{m,n\}}(\bm{A})$)
to be the largest (resp.~smallest) singular values of $\bm{A}$.

It is straightforward to read the rank of a matrix from its singular
values in that the rank is equal to the number of nonzero singular
values. If the rank $r$ of a matrix $\bm{A}\in\mathbb{R}^{m\times n}$
is smaller than $\min\{m,n\}$, it is often more economical to represent
the matrix using the \emph{compact SVD }
\begin{equation}
\bm{A}=\bm{U}_{r}\bm{\Sigma}_{r}\bm{V}_{r}^{\top},\label{eq:compact-SVD}
\end{equation}
where $\bm{U}_{r}$ (resp.~$\bm{V}_{r}$) is the first $r$ columns
of $\bm{U}$ (resp.~$\bm{V}$) and $\bm{\Sigma}_{r}$ is the upper
left block of $\bm{\Sigma}$ with size $r$.

Having introduced two decompositions, it is instrumental to compare
the singular value decomposition of a general matrix (cf.~Definition~\ref{def:SVD})
with the spectral decomposition of a real symmetric matrix (cf.~Theorem~(\ref{thm:spectral})).
First, the eigenvalues of a real symmetric matrix are not necessarily
nonnegative whereas singular values must be. Second, as claimed in
(\ref{eq:spectral-decomposition-diagonalization-form}), any real
symmetric matrix can be diagonalized by a single orthonormal matrix.
However, this may not hold true for a general matrix, nevertheless,
SVD points out a way to diagonalize a matrix using two orthonormal
matrices. To see this, one can equate (\ref{eq:SVD}) with

\[
\bm{U}^{\top}\bm{A}\bm{V}=\bm{\Sigma}.
\]
Last but not least, singular value decomposition and spectral decomposition
are inherently related. It is an easy exercise to see that the columns
of $\bm{U}$ (resp.~$\bm{V}$) are eigenvectors of $\bm{A}\bm{A}^{\top}$
(resp.~$\bm{A}^{\top}\bm{A}$). In addition, the singular values
of $\bm{A}$ are nonnegative square roots of the eigenvalues of either
$\bm{A}\bm{A}^{\top}$ or $\bm{A}^{\top}\bm{A}$ (modulo zero singular
values). This connection shall prove useful in later discussions.

\subsection{Matrix norms}

To measure the size of the perturbation (e.g.~$\bm{E}$), we need
to introduce appropriate metrics~/~norms on the space of matrices.
Among all matrix norms, unitarily invariant norms are of most interest
to us. \begin{definition}[Unitarily invariant norms]A matrix norm
$\vertiii{\cdot}$ on $\mathbb{R}^{m\times n}$ is unitarily invariant
if
\[
\vertiii{\bm{A}}=\vertiiibig{\bm{U}^{\top}\bm{A}\bm{V}}
\]
for any matrix $\bm{A}\in\mathbb{R}^{m\times n}$ and any two orthogonal
matrices $\bm{U}\in\mathbb{R}^{m\times m}$ and $\bm{V}\in\mathbb{R}^{n\times n}$.
\end{definition}

Two prominent examples include the Frobenius norm and the operator
norm

\begin{equation}
\left\Vert \bm{A}\right\Vert _{\mathrm{F}}=\sqrt{\sum_{i,j}A_{ij}^{2}}\qquad\text{and}\qquad\left\Vert \bm{A}\right\Vert =\max_{\bm{x}\in\mathbb{R}^{n}:\bm{x}\neq\bm{0}}\frac{\left\Vert \bm{A}\bm{x}\right\Vert _{2}}{\left\Vert \bm{x}\right\Vert _{2}}.\label{eq:def-fro-op-norm}
\end{equation}
To see why both $\|\cdot\|_{\mathrm{F}}$ and $\|\cdot\|$ are unitarily
invariant norms, we have the following important characterizations
\[
\left\Vert \bm{A}\right\Vert _{\mathrm{F}}=\sqrt{\sum_{i=1}^{\min\{m,n\}}\sigma_{i}^{2}\left(\bm{A}\right)}\qquad\text{and}\qquad\left\Vert \bm{A}\right\Vert _{\mathsf{2}}=\sigma_{1}\left(\bm{A}\right),
\]
where we recall that $\sigma_{1}(\bm{A})\geq\sigma_{2}(\bm{A})\geq\cdots\geq\sigma_{\min\{m,n\}}(\bm{A})\geq0$
are the singular values of $\bm{A}$ in descending order.

Last but not least, unitarily invariant norms enjoy the following
nice properties.

\begin{proposition}\label{prop:unitary_norm_relation}Let $\vertiii{\cdot}$ be any unitarily invariant norm.
One has 
\begin{align*}
\vertiii{\bm{A}\bm{B}} & \leq\vertiii{\bm{A}}\left\Vert \bm{B}\right\Vert _{2}, &  & \vertiii{\bm{A}\bm{B}}\leq\vertiii{\bm{B}}\left\Vert \bm{A}\right\Vert _{2},\\
\vertiii{\bm{A}\bm{B}} & \geq\vertiii{\bm{A}}\sigma_{\min}\left(\bm{B}\right), &  & \vertiii{\bm{A}\bm{B}}\geq\vertiii{\bm{B}}\sigma_{\min}\left(\bm{A}\right).
\end{align*}
\end{proposition}
\begin{proof}See Theorem 3.9 in \citep{stewart1990matrix}.\end{proof}

\subsection{Perturbation bounds for eigenvalues and singular values}


In this section, we review classical perturbation
bounds for eigenvalues of symmetric matrices and singular values for general matrices. 

\begin{lemma}[Weyls's inequality]\label{lemma:weyl}Let
$\bm{A},\bm{E}\in\mathbb{R}^{n\times n}$ be two real symmetric matrices.
For every $1\leq i\leq n$, we have 
\begin{equation}
\left|\lambda_{i}\left(\bm{A}\right)-\lambda_{i}\left(\bm{A} +\bm{E}\right)\right|\leq\left\Vert \bm{E}\right\Vert .\label{eq:weyl}
\end{equation}
\end{lemma}
\begin{proof}See Equation (1.63) in \citep{Tao2012RMT}. \end{proof}


An immediate implication of Lemma~\ref{lemma:weyl} is that the eigenvalues
of a real symmetric matrix are stable against perturbations. Similar bounds hold for singular values.


\begin{lemma}[Wely's inequality for singular values]\label{lemma:weyls-singular-value}Let
$\bm{A},\bm{E}\in\mathbb{R}^{m\times n}$ be two general matrices.
Without loss of generality, assume that $m\leq n$. Then for every
$1\leq i\leq m$, one has 
\[
\left|\sigma_{i}\left(\bm{A}+\bm{E}\right)-\sigma_{i}\left(\bm{A}\right)\right|\leq\left\Vert \bm{E}\right\Vert .
\]
\end{lemma}
\begin{proof}See Exercise 1.3.22 in \citep{Tao2012RMT}. \end{proof}
Again, Lemma~\ref{lemma:weyls-singular-value} reveals that singular
values are \emph{stable }with respect to small perturbations.


%\chapter{Random matrix theory}
%
%\cm{Need to motivate why we need these. }
%
%As we have mentioned in Section~XXX, we often observe
%
%One quantitative version of the central limit theorem \cm{TODO}
%
%
%\begin{theorem}[Bernstein's inequality]Let $\{x_{i}\}_{1\leq i\leq n}$
%be a sequence of independent random variables. Suppose that for each
%$1\leq i\leq n$, $\mathbb{E}[x_{i}]=0$ and $|x_{i}|\leq B$ almost
%surely. Then we have for all $t\geq0$ 
%\[
%\mathbb{P}\left(\left|\sum_{i=1}^{n}x_{i}\right|\geq t\right)\leq2\exp\left(\frac{-t^{2}/2}{V+Bt/3}\right),
%\]
%where $V$ is the variance statistic 
%\[
%V\coloneqq\left|\sum_{i=1}^{n}\mathbb{E}\left[x_{i}^{2}\right]\right|.
%\]
%\end{theorem}


